\documentclass[dvisvgm,tikz,border=3pt]{standalone}
\usepackage{amsmath,amssymb}
\usepackage{pgfplots}
\usepackage{CJKutf8}
\pgfplotsset{compat=1.18}
\usetikzlibrary{arrows.meta,calc}

\definecolor{themebg}{HTML}{2e362c}
\definecolor{themefg}{HTML}{edf4e8}
\definecolor{themegray}{HTML}{9ea897}
\definecolor{themecurve}{HTML}{7eb6ff}
\definecolor{themealert}{HTML}{ff7b7b}
\definecolor{themeamber}{HTML}{f5b942}

\begin{document}
\begin{CJK*}{UTF8}{gbsn}
\pagecolor{themebg}
\begin{tikzpicture}[color=themefg, text=themefg]

% Subplot 1: Injective (x^3)
\begin{axis}[
    name=ax1,
    width=6.0cm,
    height=6.2cm,
    axis lines=middle,
    axis line style={color=themefg, -{Stealth}},
    tick style={color=themefg},
    tick label style={color=themefg, font=\footnotesize},
    every axis label/.style={color=themefg},
    xmin=-1.9, xmax=2.1,
    ymin=-2.2, ymax=2.5,
    xtick={-1, 0, 1},
    ytick={-1, 0, 1},
    clip=false,
    xlabel={$x$},
    ylabel={$y$},
    title style={at={(0.5,1.08)},anchor=south,color=themefg,font=\bfseries\small},
    title={① 单射型：每个水平线至多一交点},
]
  % Curve y = x^3
  \addplot[themecurve, line width=2pt, domain=-1.3:1.3, samples=150] {x^3};
  \node[text=themecurve, fill=themebg, inner sep=0.8pt, anchor=south east] at (axis cs:1.35,1.4) {\footnotesize $y=x^3$};

  % Horizontal test line y = 0.8
  \draw[dashed, themegray, thick] (axis cs:-1.7,0.8) -- (axis cs:1.7,0.8);
  \node[text=themegray, fill=themebg, inner sep=0.8pt, anchor=south east] at (axis cs:-0.9,0.82) {\scriptsize $y_0$};

  % Single intersection
  \addplot[only marks, mark=*, mark size=2.5pt, fill=themecurve, draw=themefg] coordinates {(0.9283, 0.8)};
  \draw[dotted, themegray, thick] (axis cs:0.9283, 0.8) -- (axis cs:0.9283, 0);
  \node[text=themefg, fill=themebg, inner sep=0.8pt, below] at (axis cs:0.9283, -0.05) {\footnotesize $x_0$};
\end{axis}

\node[font=\small\bfseries, color=themefg, anchor=north] at ($(ax1.south)-(0, 0.45cm)$) {原型：奇次幂 / 指对数};

% Subplot 2: Folding (x^2)
\begin{axis}[
    at={($(ax1.east)+(2.0cm,0)$)},
    anchor=west,
    name=ax2,
    width=6.0cm,
    height=6.2cm,
    axis lines=middle,
    axis line style={color=themefg, -{Stealth}},
    tick style={color=themefg},
    tick label style={color=themefg, font=\footnotesize},
    every axis label/.style={color=themefg},
    xmin=-1.9, xmax=2.1,
    ymin=-0.6, ymax=2.9,
    xtick={-1, 0, 1},
    ytick={0, 1, 2},
    clip=false,
    xlabel={$x$},
    ylabel={$y$},
    title style={at={(0.5,1.08)},anchor=south,color=themefg,font=\bfseries\small},
    title={② 二原像示例：$y=x^2$ 在 $y_0>0$ 时有两个原像},
]
  % Curve y = x^2
  \addplot[themealert, line width=2pt, domain=-1.55:1.55, samples=150] {x^2};
  \node[text=themealert, fill=themebg, inner sep=0.8pt, anchor=south east] at (axis cs:1.6,1.8) {\footnotesize $y=x^2$};

  % Horizontal test line y = 1.0
  \draw[dashed, themegray, thick] (axis cs:-1.7,1.0) -- (axis cs:1.7,1.0);
  \node[text=themegray, fill=themebg, inner sep=0.8pt, anchor=south east] at (axis cs:-1.1,1.02) {\scriptsize $y_0$};

  % Dual intersections
  \addplot[only marks, mark=*, mark size=2.5pt, fill=themealert, draw=themefg] coordinates {(-1.0, 1.0) (1.0, 1.0)};
  \draw[dotted, themegray, thick] (axis cs:-1.0, 1.0) -- (axis cs:-1.0, 0);
  \draw[dotted, themegray, thick] (axis cs:1.0, 1.0) -- (axis cs:1.0, 0);
  \node[text=themefg, fill=themebg, inner sep=0.8pt, below] at (axis cs:-1.0, -0.05) {\footnotesize $-x_0$};
  \node[text=themefg, fill=themebg, inner sep=0.8pt, below] at (axis cs:1.0, -0.05) {\footnotesize $x_0$};
\end{axis}

\node[font=\small\bfseries, color=themefg, anchor=north] at ($(ax2.south)-(0, 0.45cm)$) {原型：偶次幂 / $\lvert x \rvert$};

% Subplot 3: Periodic (sin x)
\begin{axis}[
    at={($(ax1.south)+(4.0cm,-2.0cm)$)},
    anchor=north,
    name=ax3,
    width=6.0cm,
    height=6.2cm,
    axis lines=middle,
    axis line style={color=themefg, -{Stealth}},
    tick style={color=themefg},
    tick label style={color=themefg, font=\footnotesize},
    every axis label/.style={color=themefg},
    xmin=-3.8, xmax=4.1,
    ymin=-1.6, ymax=1.9,
    xtick={-3.14159, 0, 3.14159},
    xticklabels={$-\pi$, $0$, $\pi$},
    ytick={-1, 0, 1},
    clip=false,
    xlabel={$x$},
    ylabel={$y$},
    title style={at={(0.5,1.08)},anchor=south,color=themefg,font=\bfseries\small},
    title={③ 周期重复型：$y=\sin x$ 的同一函数值周期性重复},
]
  % Curve y = sin x
  \addplot[themecurve, line width=2pt, domain=-3.75:3.75, samples=220] {sin(deg(x))};
  \node[text=themecurve, fill=themebg, inner sep=0.8pt, anchor=south east] at (axis cs:3.5,0.7) {\footnotesize $y=\sin x$};

  % Horizontal test line y = 0.5
  \draw[dashed, themegray, thick] (axis cs:-3.7,0.5) -- (axis cs:3.7,0.5);
  \node[text=themegray, fill=themebg, inner sep=0.8pt, anchor=south east] at (axis cs:-2.8,0.52) {\scriptsize $y_0$};

  % Multiple intersections
  \addplot[only marks, mark=*, mark size=2.5pt, fill=themecurve, draw=themefg] coordinates {(-3.6652, 0.5) (0.5236, 0.5) (2.618, 0.5)};
  \draw[dotted, themegray, thick] (axis cs:0.5236, 0.5) -- (axis cs:0.5236, 0);
  \draw[dotted, themegray, thick] (axis cs:2.618, 0.5) -- (axis cs:2.618, 0);
\end{axis}

\node[font=\small\bfseries, color=themefg, anchor=north] at ($(ax3.south)-(0, 0.45cm)$) {原型：三角函数周期律};

\end{tikzpicture}
\end{CJK*}
\end{document}
